% Generated draft: requires mathematical and source review.
\documentclass[11pt]{article}
\usepackage[margin=25mm]{geometry}
\usepackage{fontspec}
\setmainfont{Latin Modern Roman}
\usepackage{amsmath,amssymb,mathtools,unicode-math}
\setmathfont{Latin Modern Math}
\usepackage{enumitem,needspace,xurl,hyperref}
\setlength{\parindent}{0pt}
\setlength{\parskip}{5pt}
\setlength{\emergencystretch}{3em}
\newcommand{\N}{\mathbb{N}}
\newcommand{\Z}{\mathbb{Z}}
\newcommand{\Q}{\mathbb{Q}}
\newcommand{\R}{\mathbb{R}}
\newcommand{\C}{\mathbb{C}}
\newcommand{\F}{\mathbb{F}}
\newcommand{\A}{\mathbb{A}}
\newcommand{\PP}{\mathbb{P}}
\newcommand{\E}{\mathbb{E}}
\newcommand{\eps}{\varepsilon}
\newcommand{\dd}{\,\mathrm{d}}
\begin{document}
\section*{UnsolvedMath ID: 11000067}
\textbf{Source problem code:} catalogue entry
\subsection{Definitions and mathematical statement}
In the usual closed oriented genus-\(g\) surface convention, the level-\(L\) subgroup is
\[
 \operatorname{Mod}_g[L]
 =\ker\bigl(\operatorname{Mod}_g\longrightarrow
       \operatorname{Sp}(H_1(\Sigma_g;\Z/L\Z))\bigr),
\]
for an allowed positive integral level \(L\). The problem is to compute its integral first homology, equivalently its abelianization,
\[
 H_1(\operatorname{Mod}_g[L];\Z)
 \cong\operatorname{Mod}_g[L]/[\operatorname{Mod}_g[L],\operatorname{Mod}_g[L]],
\]
with an explicit group description, for example its invariant factors. The exact genus and level ranges are not repeated in the problem sentence.

The source reports that for \(g\geq3\) this group is finite, proved independently by McCarthy and Hain, referring to [Ha2, Proposition 5.2]. Thus computing it asks for more than establishing finiteness. A footnote says Hain proves a stronger result computing \(H^1(\operatorname{Mod}_g[L];V)\) for finite-dimensional symplectic representations \(V\); that is a cohomology statement with coefficients and must not be confused with the requested integral homology group.

The continuation says a following conjecture, discussed in [Ha2, Section 5], would imply rank one for the orbifold Picard group of the moduli space of level-\(L\) structures, and reports that Picard group's finite generation from the preceding Hain--McCarthy result. The following conjecture is absent. These are reported motivations, not an asserted rank-one consequence of the computation alone.
\subsection{Short English statement}
Compute the full abelianization of the level mapping class subgroup, beyond its reported finiteness for \(g\geq3\), retaining the distinction from Hain's coefficient-cohomology result.
\subsection{Sources}
\begin{itemize}
\item \emph{Problems on Mapping Class Groups and Related Topics}, edited by Farb, Problem 5.23 and its continuation/footnote, supplied PDF page 51. \url{https://www.math.uchicago.edu/~farb/papers/mcgbook.pdf}. Recorded access 2026-07-29; extraction pdf-font-structure.
\item McCarthy and Hain, independent finiteness attribution; Hain [Ha2, Proposition 5.2 and Section 5]. The complete bibliography entry is missing. The footnote credits Hain's stronger coefficient-cohomology computation.
\item Local UnsolvedMath registry, \texttt{data/unsolvedmath/problems.json}, permanent ID \(11000067\), exact source code \texttt{catalogue entry}. Its OPEN-TRIAGE summary mentions a level-\(2\) partial result with truncated attribution ``Sur-...''; no complete author or citation is supplied.
\end{itemize}
\end{document}
